\iftestbuild % TEST-ONLY-BEGIN
\setcounter{section}{2}\setcounter{theorem}{33}
\fi % TEST-ONLY-END
% Source: book pp. 251--257 (PDF 265--271); no solutions.
% \axeqno: an equation number in Axler's item counter, inside \[ ... \]:
%   \[ ... \axeqno\label{ax:7.42} \]
\DeclareRobustCommand\axeqno{\refstepcounter{theorem}\tag*{\thetheorem}}
\section{Positive Operators}

\begin{definition}[Positive operator]\label{ax:7.34}
An operator $T\in\Lin(V)$ is called \emph{positive} if $T$ is self-adjoint and
\[ \ip{Tv}{v}\ge0 \]
for all $v\in V$.
\end{definition}

If $V$ is a complex vector space, then the requirement that $T$ be self-adjoint
can be dropped from the definition above (by \ref{ax:7.14}).

\begin{example}[Positive operators]\label{ax:7.35}\leavevmode
\begin{parts}
\item Let $T\in\Lin(\F^2)$ be the operator whose matrix (using the standard
  basis) is $\bigl(\begin{smallmatrix}2&-1\\-1&1\end{smallmatrix}\bigr)$. Then
  $T$ is self-adjoint and
  $\ip{T(w,z)}{(w,z)}=2\abs{w}^2-2\Re(w\overline{z})+\abs{z}^2
  =\abs{w-z}^2+\abs{w}^2\ge0$ for all $(w,z)\in\F^2$. Thus $T$ is a positive
  operator.
\item If $U$ is a subspace of $V$, then the orthogonal projection $P_U$ is a
  positive operator, as you should verify.
\item If $T\in\Lin(V)$ is self-adjoint and $b,c\in\R$ are such that $b^2<4c$,
  then $T^2+bT+cI$ is a positive operator, as shown by the proof of
  \ref{ax:7.26}.
\end{parts}
\end{example}

\begin{definition}[Square root]\label{ax:7.36}
An operator $R$ is called a \emph{square root} of an operator $T$ if $R^2=T$.
\end{definition}

\begin{example}[Square root of an operator]\label{ax:7.37}
If $T\in\Lin(\F^3)$ is defined by $T(z_1,z_2,z_3)=(z_3,0,0)$, then the operator
$R\in\Lin(\F^3)$ defined by $R(z_1,z_2,z_3)=(z_2,z_3,0)$ is a square root of $T$
because $R^2=T$, as you can verify.
\end{example}

\marginnote{Because positive operators correspond to nonnegative numbers, better
terminology would use the term \emph{nonnegative operators}. However, operator
theorists consistently call these positive operators, so we follow that custom.
Some mathematicians use the term \textbf{positive semidefinite operator}, which
means the same as positive operator.}%
The characterizations of the positive operators in the next result correspond
to characterizations of the nonnegative numbers among $\C$. Specifically, a
number $z\in\C$ is nonnegative if and only if it has a nonnegative square root,
corresponding to condition (d). Also, $z$ is nonnegative if and only if it has a
real square root, corresponding to condition (e). Finally, $z$ is nonnegative if
and only if there exists $w\in\C$ such that $z=\overline{w}w$, corresponding to
condition (f). See Exercise~20 for another condition that is equivalent to
being a positive operator.

\begin{result}[Characterizations of positive operators]\label{ax:7.38}
Let $T\in\Lin(V)$. Then the following are equivalent.
\begin{parts}
\item $T$ is a positive operator.
\item $T$ is self-adjoint and all eigenvalues of $T$ are nonnegative.
\item With respect to some orthonormal basis of $V$, the matrix of $T$ is a
  diagonal matrix with only nonnegative numbers on the diagonal.
\item $T$ has a positive square root.
\item $T$ has a self-adjoint square root.
\item $T=R^*R$ for some $R\in\Lin(V)$.
\end{parts}
\end{result}

\begin{proof}
We will prove that (a) $\Rightarrow$ (b) $\Rightarrow$ (c) $\Rightarrow$ (d)
$\Rightarrow$ (e) $\Rightarrow$ (f) $\Rightarrow$ (a).

First suppose (a) holds, so that $T$ is positive, which implies that $T$ is
self-adjoint (by definition of positive operator). To prove the other condition
in (b), suppose $\lambda$ is an eigenvalue of $T$. Let $v$ be an eigenvector of
$T$ corresponding to $\lambda$. Then
\[ 0\le\ip{Tv}{v}=\ip{\lambda v}{v}=\lambda\ip{v}{v}. \]
Thus $\lambda$ is a nonnegative number. Hence (b) holds, showing that (a)
implies (b).

Now suppose (b) holds, so that $T$ is self-adjoint and all eigenvalues of $T$
are nonnegative. By the spectral theorem (\ref{ax:7.29} and \ref{ax:7.31}),
there is an orthonormal basis $e_1,\dots,e_n$ of $V$ consisting of eigenvectors
of $T$. Let $\lambda_1,\dots,\lambda_n$ be the eigenvalues of $T$ corresponding
to $e_1,\dots,e_n$; thus each $\lambda_k$ is a nonnegative number. The matrix of
$T$ with respect to $e_1,\dots,e_n$ is the diagonal matrix with
$\lambda_1,\dots,\lambda_n$ on the diagonal, which shows that (b) implies (c).

Now suppose (c) holds. Suppose $e_1,\dots,e_n$ is an orthonormal basis of $V$
such that the matrix of $T$ with respect to this basis is a diagonal matrix with
nonnegative numbers $\lambda_1,\dots,\lambda_n$ on the diagonal. The linear map
lemma (\ref{ax:3.4}) implies that there exists $R\in\Lin(V)$ such that
\[ Re_k=\sqrt{\lambda_k}\,e_k \]
for each $k=1,\dots,n$. As you should verify, $R$ is a positive operator.
Furthermore, $R^2e_k=\lambda_ke_k=Te_k$ for each $k$, which implies that
$R^2=T$. Thus $R$ is a positive square root of $T$. Hence (d) holds, which shows
that (c) implies (d).

Every positive operator is self-adjoint (by definition of positive operator).
Thus (d) implies (e).

Now suppose (e) holds, meaning that there exists a self-adjoint operator $R$ on
$V$ such that $T=R^2$. Then $T=R^*R$ (because $R^*=R$). Hence (e) implies (f).

Finally, suppose (f) holds. Let $R\in\Lin(V)$ be such that $T=R^*R$. Then
$T^*=(R^*R)^*=R^*(R^*)^*=R^*R=T$. Hence $T$ is self-adjoint. To complete the
proof that (a) holds, note that
\[ \ip{Tv}{v}=\ip{R^*Rv}{v}=\ip{Rv}{Rv}\ge0 \]
for every $v\in V$. Thus $T$ is positive, showing that (f) implies (a).
\end{proof}

Every nonnegative number has a unique nonnegative square root. The next result
shows that positive operators enjoy a similar property.

\begin{result}[Each positive operator has only one positive square root]\label{ax:7.39}
Every positive operator on $V$ has a unique positive square root.
\end{result}

\begin{proof}
Suppose $T\in\Lin(V)$ is positive. Suppose $v\in V$ is an eigenvector of $T$.
Hence there exists a real number $\lambda\ge0$ such that $Tv=\lambda v$.

Let $R$ be a positive square root of $T$. We will prove that
$Rv=\sqrt{\lambda}\,v$. This will imply that the behavior of $R$ on the
eigenvectors of $T$ is uniquely determined. Because there is a basis of $V$
consisting of eigenvectors of $T$ (by the spectral theorem), this will imply
that $R$ is uniquely determined.

\marginnote{A positive operator can have infinitely many square roots (although
only one of them can be positive). For example, the identity operator on $V$
has infinitely many square roots if $\dim V>1$.}%

To prove that $Rv=\sqrt{\lambda}\,v$, note that the spectral theorem asserts that
there is an orthonormal basis $e_1,\dots,e_n$ of $V$ consisting of eigenvectors
of $R$. Because $R$ is a positive operator, all its eigenvalues are nonnegative.
Thus there exist nonnegative numbers $\lambda_1,\dots,\lambda_n$ such that
$Re_k=\sqrt{\lambda_k}\,e_k$ for each $k=1,\dots,n$.

Because $e_1,\dots,e_n$ is a basis of $V$, we can write
\[ v=a_1e_1+\cdots+a_ne_n \]
for some numbers $a_1,\dots,a_n\in\F$. Thus
\[ Rv=a_1\sqrt{\lambda_1}\,e_1+\cdots+a_n\sqrt{\lambda_n}\,e_n. \]
Hence
\[ \lambda v=Tv=R^2v=a_1\lambda_1e_1+\cdots+a_n\lambda_ne_n. \]
The equation above implies that
\[ a_1\lambda e_1+\cdots+a_n\lambda e_n=a_1\lambda_1e_1+\cdots+a_n\lambda_ne_n. \]
Thus $a_k(\lambda-\lambda_k)=0$ for each $k=1,\dots,n$. Hence
\[ v=\sum_{\{k\,:\,\lambda_k=\lambda\}}a_ke_k. \]
Thus
\[ Rv=\sum_{\{k\,:\,\lambda_k=\lambda\}}a_k\sqrt{\lambda}\,e_k=\sqrt{\lambda}\,v, \]
as desired.
\end{proof}

The notation defined below makes sense thanks to the result above.

\begin{notation}[$\sqrt{T}$]\label{ax:7.40}
For $T$ a positive operator, $\sqrt{T}$ denotes the unique positive square root
of $T$.
\end{notation}

\begin{example}[Square root of positive operators]\label{ax:7.41}
Define operators $S,T$ on $\R^2$ (with the usual Euclidean inner product) by
\[ S(x,y)=(x,2y) \quad\text{and}\quad T(x,y)=(x+y,x+y). \]
Then with respect to the standard basis of $\R^2$ we have
\[ \Mat(S)=\begin{pmatrix}1&0\\0&2\end{pmatrix} \quad\text{and}\quad
   \Mat(T)=\begin{pmatrix}1&1\\1&1\end{pmatrix}. \axeqno\label{ax:7.42} \]
Each of these matrices equals its transpose; thus $S$ and $T$ are self-adjoint.

If $(x,y)\in\R^2$, then
\[ \ip{S(x,y)}{(x,y)}=x^2+2y^2\ge0 \]
and
\[ \ip{T(x,y)}{(x,y)}=x^2+2xy+y^2=(x+y)^2\ge0. \]
Thus $S$ and $T$ are positive operators.

The standard basis of $\R^2$ is an orthonormal basis consisting of eigenvectors
of $S$. Note that
\[ \Bigl(1/\sqrt2,1/\sqrt2\Bigr),
   \Bigl(1/\sqrt2,-1/\sqrt2\Bigr) \]
is an orthonormal basis of eigenvectors of $T$, with eigenvalue $2$ for the
first eigenvector and eigenvalue $0$ for the second eigenvector. Thus $\sqrt{T}$
has the same eigenvectors, with eigenvalues $\sqrt2$ and $0$.

You can verify that
\[ \Mat\bigl(\sqrt{S}\bigr)=\begin{pmatrix}1&0\\0&\sqrt2\end{pmatrix}
   \quad\text{and}\quad
   \Mat\bigl(\sqrt{T}\bigr)=\begin{pmatrix}1/\sqrt2&1/\sqrt2\\[2pt]
   1/\sqrt2&1/\sqrt2\end{pmatrix} \]
with respect to the standard basis by showing that the squares of the matrices
above are the matrices in \ref{ax:7.42} and that each matrix above is the matrix
of a positive operator.
\end{example}

The statement of the next result does not involve a square root, but the clean
proof makes nice use of the square root of a positive operator.

\begin{result}[$T$ positive and $\ip{Tv}{v}=0\implies Tv=0$]\label{ax:7.43}
Suppose $T$ is a positive operator on $V$ and $v\in V$ is such that
$\ip{Tv}{v}=0$. Then $Tv=0$.
\end{result}

\begin{proof}
We have
\[ 0=\ip{Tv}{v}=\bigl\langle\sqrt{T}\sqrt{T}v,v\bigr\rangle
   =\bigl\langle\sqrt{T}v,\sqrt{T}v\bigr\rangle=\bigl\lVert\sqrt{T}v\bigr\rVert^2. \]
Hence $\sqrt{T}v=0$. Thus $Tv=\sqrt{T}\bigl(\sqrt{T}v\bigr)=0$, as desired.
\end{proof}

% ------------------------------------------------------------------ exercises
\exercisesheading{7C}

\begin{exercise}
Suppose $T\in\Lin(V)$. Prove that if both $T$ and $-T$ are positive operators,
then $T=0$.
\end{exercise}

\begin{exercise}
Suppose $T\in\Lin(\F^4)$ is the operator whose matrix (with respect to the
standard basis) is
\[ \begin{pmatrix}2&-1&0&0\\-1&2&-1&0\\0&-1&2&-1\\0&0&-1&2\end{pmatrix}. \]
Show that $T$ is an invertible positive operator.
\end{exercise}

\begin{exercise}
Suppose $n$ is a positive integer and $T\in\Lin(\F^n)$ is the operator whose
matrix (with respect to the standard basis) consists of all $1$'s. Show that $T$
is a positive operator.
\end{exercise}

\begin{exercise}
Suppose $n$ is an integer with $n>1$. Show that there exists an $n$-by-$n$
matrix $A$ such that all of the entries of $A$ are positive numbers and
$A=A^*$, but the operator on $\F^n$ whose matrix (with respect to the standard
basis) equals $A$ is not a positive operator.
\end{exercise}

\begin{exercise}
Suppose $T\in\Lin(V)$ is self-adjoint. Prove that $T$ is a positive operator if
and only if for every orthonormal basis $e_1,\dots,e_n$ of $V$, all entries on
the diagonal of $\Mat\bigl(T,(e_1,\dots,e_n)\bigr)$ are nonnegative numbers.
\end{exercise}

\begin{exercise}
Prove that the sum of two positive operators on $V$ is a positive operator.
\end{exercise}

\begin{exercise}
Suppose $S\in\Lin(V)$ is an invertible positive operator and $T\in\Lin(V)$ is a
positive operator. Prove that $S+T$ is invertible.
\end{exercise}

\begin{exercise}
Suppose $T\in\Lin(V)$. Prove that $T$ is a positive operator if and only if the
pseudoinverse $T^\dagger$ is a positive operator.
\end{exercise}

\begin{exercise}
Suppose $T\in\Lin(V)$ is a positive operator and $S\in\Lin(W,V)$. Prove that
$S^*TS$ is a positive operator on $W$.
\end{exercise}

\begin{exercise}
Suppose $T$ is a positive operator on $V$. Suppose $v,w\in V$ are such that
\[ Tv=w \quad\text{and}\quad Tw=v. \]
Prove that $v=w$.
\end{exercise}

\begin{exercise}
Suppose $T$ is a positive operator on $V$ and $U$ is a subspace of $V$
invariant under $T$. Prove that $T|_U\in\Lin(U)$ is a positive operator on $U$.
\end{exercise}

\begin{exercise}
Suppose $T\in\Lin(V)$ is a positive operator. Prove that $T^k$ is a positive
operator for every positive integer $k$.
\end{exercise}

\begin{exercise}
Suppose $T\in\Lin(V)$ is self-adjoint and $\alpha\in\R$.
\begin{parts}
\item Prove that $T-\alpha I$ is a positive operator if and only if $\alpha$ is
  less than or equal to every eigenvalue of $T$.
\item Prove that $\alpha I-T$ is a positive operator if and only if $\alpha$ is
  greater than or equal to every eigenvalue of $T$.
\end{parts}
\end{exercise}

\begin{exercise}
Suppose $T$ is a positive operator on $V$ and $v_1,\dots,v_m\in V$. Prove that
\[ \sum_{j=1}^m\sum_{k=1}^m\ip{Tv_k}{v_j}\ge0. \]
\end{exercise}

\begin{exercise}
Suppose $T\in\Lin(V)$ is self-adjoint. Prove that there exist positive operators
$A,B\in\Lin(V)$ such that
\[ T=A-B \quad\text{and}\quad \sqrt{T^*T}=A+B \quad\text{and}\quad AB=BA=0. \]
\end{exercise}

\begin{exercise}
Suppose $T$ is a positive operator on $V$. Prove that
\[ \nul\sqrt{T}=\nul T \quad\text{and}\quad \range\sqrt{T}=\range T. \]
\end{exercise}

\begin{exercise}
Suppose that $T\in\Lin(V)$ is a positive operator. Prove that there exists a
polynomial $p$ with real coefficients such that $\sqrt{T}=p(T)$.
\end{exercise}

\begin{exercise}
Suppose $S$ and $T$ are positive operators on $V$. Prove that $ST$ is a positive
operator if and only if $S$ and $T$ commute.
\end{exercise}

\begin{exercise}
Show that the identity operator on $\F^2$ has infinitely many self-adjoint
square roots.
\end{exercise}

\begin{exercise}
Suppose $T\in\Lin(V)$ and $e_1,\dots,e_n$ is an orthonormal basis of $V$. Prove
that $T$ is a positive operator if and only if there exist $v_1,\dots,v_n\in V$
such that
\[ \ip{Te_k}{e_j}=\ip{v_k}{v_j} \]
for all $j,k=1,\dots,n$.
\exnote{The numbers $\{\ip{Te_k}{e_j}\}_{j,k=1,\dots,n}$ are the entries in
the matrix of $T$ with respect to the orthonormal basis $e_1,\dots,e_n$.}
\end{exercise}

\begin{exercise}
Suppose $n$ is a positive integer. The $n$-by-$n$ \emph{Hilbert matrix} is the
$n$-by-$n$ matrix whose entry in row $j$, column $k$ is $1/(j+k-1)$.
Suppose $T\in\Lin(V)$ is an operator whose matrix with respect to some
orthonormal basis of $V$ is the $n$-by-$n$ Hilbert matrix. Prove that $T$ is a
positive invertible operator.
\exnote{Example: The $4$-by-$4$ Hilbert matrix is
\[ \begin{pmatrix}
   1&1/2&1/3&1/4\\[3pt]
   1/2&1/3&1/4&1/5\\[3pt]
   1/3&1/4&1/5&1/6\\[3pt]
   1/4&1/5&1/6&1/7
   \end{pmatrix}. \]}
\end{exercise}

\begin{exercise}
Suppose $T\in\Lin(V)$ is a positive operator and $u\in V$ is such that
$\norm{u}=1$ and $\norm{Tu}\ge\norm{Tv}$ for all $v\in V$ with $\norm{v}=1$.
Show that $u$ is an eigenvector of $T$ corresponding to the largest eigenvalue
of $T$.
\end{exercise}

\begin{exercise}
For $T\in\Lin(V)$ and $u,v\in V$, define $\ip{u}{v}_T$ by
$\ip{u}{v}_T=\ip{Tu}{v}$.
\begin{parts}
\item Suppose $T\in\Lin(V)$. Prove that $\ip{\cdot}{\cdot}_T$ is an inner
  product on $V$ if and only if $T$ is an invertible positive operator (with
  respect to the original inner product $\ip{\cdot}{\cdot}$).
\item Prove that every inner product on $V$ is of the form
  $\ip{\cdot}{\cdot}_T$ for some positive invertible operator $T\in\Lin(V)$.
\end{parts}
\end{exercise}

\begin{exercise}
Suppose $S$ and $T$ are positive operators on $V$. Prove that
\[ \nul(S+T)=\nul S\cap\nul T. \]
\end{exercise}

\begin{exercise}
Let $T$ be the second derivative operator in Exercise~31(b) in Section~7A. Show
that $-T$ is a positive operator.
\end{exercise}
